\section*{Resumo}

Este trabalho investiga a aplicação da abordagem retrieval-augmented generation (RAG) no desenvolvimento de uma prova de conceito de chatbot legislativo voltada à consulta de discursos da 56\textordfeminine{} Legislatura do Senado Federal (2019--2023). Parte-se do problema da dificuldade de acesso e recuperação de informação sobre discursos parlamentares por meio de consultas em linguagem natural, especialmente quando se exige precisão factual, contextualização e rastreabilidade das informações. A metodologia consistiu na construção de um pipeline de processamento com extração e preparação de dados, segmentação textual (chunking), enriquecimento por metadados, geração de embeddings, indexação vetorial e implementação de fluxo de consulta baseado em LLM. A solução foi desenvolvida em ambiente reprodutível apoiado por Open WebUI, ChromaDB, API da OpenAI na configuração experimental validada e scripts próprios para ingestão e avaliação. A demonstração empírica combinou inspeção manual, três rodadas automatizadas com rubrica e avaliação por LLM como juiz, além de análise manual amostral. Os resultados indicam bom desempenho em perguntas factuais, temáticas e de autoria bem delimitada, com médias entre 9,15 e 9,55 nas rodadas automatizadas. Observou-se, contudo, maior instabilidade em consultas que exigem precisão numérica, comparação entre autores, controle de escopo ou síntese multifocal. Conclui-se que os resultados oferecem evidência inicial de viabilidade da abordagem RAG para recuperação auditável de discursos legislativos, desde que acompanhada de governança de fontes, avaliação sistemática e supervisão humana em respostas de maior abstração.

\textbf{Palavras-chave:} RAG; inteligência artificial; chatbot legislativo; recuperação de informação; Senado Federal; LLM.
